\scanpage{54}
\mspage{10}
% archive sheet number: 53
des graphes pondérés à repères, l'opération \deleted{s'explicite} se reflète par une opération correspondante au niveau de $\vec R$, par le ``gommage'' d'une partie $\vec R'$ stable par $\underline{\unsure{D|\mathbb D}}_\infty$ i.e. par $\sigma_0,\sigma_1$

2) Gommage de cercles de découpage

On se donne maintenant \deleted{une ens. de composantes connexes de} \unsure{1-|l-}sous-variété \inserted{\unsure{(ouv. et)|(une fl.)} fermée} $\Sigma'$ de $\Sigma\setminus\partial X$, correspondant à un ensemble \deleted{$\unsure{A'_0}$} de comp. conn. de $\Sigma'$, i.e. à une partie \inserted{$A'_0$} de $\pi_0(\Sigma')=A_0$ (arêtes ordinaires du graphe) --- et on remplace $\Sigma$ par $\Sigma\setminus\Sigma'$. Au niveau des graphes pondérés, cela se traduit par une opération de \uline{contraction} \uline{en un point} des arêtes correspondantes, d'où formation d'un nouveau graphe, et d'un \uline{morphisme} \uline{de généralisation} (à l'opposé de ``spécialisation'') de l'ancien vers le nouveau.

Il convient ici d'expliciter comment la pondération suit\ldots\ Au niveau des \deleted{\unsure{repères}} \inserted{sommets}, il est entendu que les sommets repères se trouvant sur $\Sigma'$ sont gommés en $=$ temps --- donc, sur le graphe pondéré à repères, on fait sauter l'ens des repères qui sont au dessus des arêtes $\in A'_0$.
